\subsection{Relation to existing approaches}
  \label{subsec:relation-to-existing-approaches}

  The present construction shares motivations with
  background-independent approaches such as causal set theory,
  loop quantum gravity, and spectral geometry.
  Unlike causal sets, no fundamental discreteness is postulated;
  effective granularity arises from spectral filtering.
  In contrast to loop-based models, no fundamental spin networks
  or combinatorial structures are introduced.
  Compared to noncommutative geometry,
  the focus is not on algebraic generalization of manifolds
  but on operational reconstruction of metric structure
  from relational correlations.

  \paragraph{Graph Laplacian convergence.}

    The reconstruction relies on established convergence results:
    graph Laplacians converge to the Laplace--Beltrami operator
    in the dense sampling limit under mild regularity assumptions~\cite{Belkin2003,Hein2007}.
    This justifies interpreting the admissible spectral sector
    as an effective continuum geometry.

    Spectral proximity is closely related to diffusion maps~\cite{Coifman2006}.
    The distance $d_S(i,j)$ corresponds to a fixed-time diffusion
    distance with a specific spectral weighting.
    The present framework differs in two respects:
    (i) spectral distance feeds a hierarchical reconstruction
    (Appendix~B) rather than serving as an embedding tool,
    and (ii) the admissibility filter $F_{\lambda^*}$
    imposes additional spectral regularity constraints.

    Varadhan’s formula
    \begin{equation}
      d(x,y)^2 = -4t \log p_t(x,y) + o(t)
      \quad \text{as } t \to 0^+,
    \end{equation}
    establishes that short-time heat-kernel data encode
    Riemannian geodesic distance~\cite{Varadhan1967}.
    The present construction exploits this correspondence
    in finite-resolution discrete form.

    The approach differs from Connes’ noncommutative geometry~\cite{Connes1994}, which reconstructs metrics from
    spectral triples $(\mathcal A,\mathcal H,D)$.
    Here only the scalar Laplacian is used,
    without invoking spinor structure or operator-algebraic axioms.
    Both approaches recover geometry from spectral data,
    but at different levels of generality.
